\newpage
\heading{19  Adquisición y contrato de datos}{IMPLEMENTACIÓN VERIFICABLE}
\textbf{Hardware contemplado por el software.} AD8232 para señal ECG y frecuencia estimada; MAX30102 para SpO\textsubscript{2} y pulso óptico; MPU6050 para estimación respiratoria por movimiento. Estos nombres describen el contrato implementado, no acreditan cableado, calibración ni firmware efectivamente instalado.
\begin{infobox}
\textbf{Ruta web:} ESP32 $\rightarrow$ notificaciones BLE $\rightarrow$ navegador $\rightarrow$ API Next.js $\rightarrow$ PostgreSQL/Supabase.\par
\textbf{Alternativa local:} ESP32 $\rightarrow$ Python/Bleak $\rightarrow$ API. Su cola SQLite conserva reintentos locales y no reemplaza la base del backend.
\end{infobox}
\textbf{Contrato CSV-11.} Cada línea contiene once campos numéricos, en este orden:
\begin{quote}\small\ttfamily ecg, bpm\_ecg, estado\_ecg\_x50, resp\_mg, rpm,\\ estado\_resp\_x50, spo2, bpm\_oxi, estado\_oxi\_x50, ir, rojo\end{quote}
El navegador reconstruye fragmentos hasta un salto de línea. Rechaza columnas vacías, números no finitos, estados no admitidos, UTF-8 inválido y acumulaciones superiores a 16 KiB. El modo de dos copias comprueba igualdad de fragmentos; se activa solo si el firmware realmente los duplica.
\heading{19.1  Calidad, tiempo y envío}{LÍMITES DEL TRANSPORTE}
Se envía como máximo una instantánea por segundo. Se deriva una medición de pulso cuando el estado ECG es 100 y el valor cumple el formato de 1 a 300 lpm. Sin contacto o sin señal se conserva telemetría y estado sin inventar un pulso de cero. Los rangos de formato no son intervalos fisiológicos de normalidad.

La fecha de captura web procede del reloj receptor, no de un reloj ADC sincronizado. La cola web conserva hasta diez paquetes pendientes en memoria, con un envío simultáneo y hasta tres intentos por paquete. Se mantienen los identificadores al reintentar; el backend detecta duplicados y conflictos de contenido. Cerrar o recargar puede perder pendientes no confirmados.

El monitor ECG web conserva hasta diez segundos y 3000 muestras antes del muestreo de la API. La escala utiliza unidades crudas, no milivoltios calibrados. El filtro es visual: no determina latidos ni altera registros guardados. Las instantáneas del backend no reconstruyen la onda de alta frecuencia.
\vfill
{\footnotesize\color{muted}Evidencia: \path{features/capture/protocol.ts}, \path{BleCapture.tsx}, \path{ecgStream.ts}; \path{docs/CAPTURA_BLE_WEB.md}, \path{docs/PUENTE_ESP32.md}. Rutas de frontend relativas a \path{src/frontend/src/frontend/}.}
\newpage
\heading{20  Método de evaluación}{DOS ANÁLISIS COMPLEMENTARIOS}
\textbf{Referencia personal.} Se comparan mediciones de la misma persona, misión, indicador, unidad, actividad y revisión ambiental. Se separan registros simulados de no simulados y se eliminan IDs repetidos. Para pulso se exige reposo confirmado. Fecha futura, mala calidad o caducidad impiden una evaluación vigente.
\par\noindent\begin{tabularx}{\linewidth}{@{}lXXX@{}}
\toprule Parámetro & Pulso en reposo & Sueño & Bienestar\\\midrule
Ventana histórica & 14 días & 14 días & 14 días\\
Referencia mínima & 5 muestras & 5 muestras & 5 muestras\\
Lecturas consecutivas & 3 & 3 & 3\\
Desviación absoluta & 15 lpm & 1,5 h & 1 punto\\
Vigencia máxima & 30 s & 2 días & 2 días\\\bottomrule
\end{tabularx}\par
La media de referencia excluye el bloque de tres lecturas evaluado. Para cada una, $\Delta_i=x_i-\bar{x}_{ref}$. Hay desviación sostenida si las tres diferencias alcanzan el umbral positivo o las tres son como máximo su negativo. Sin referencia suficiente se informa insuficiencia; no se afirma estabilidad. Versiones: \texttt{demo-heart-rate-rest-2}, \texttt{demo-sleep-1}, \texttt{demo-wellbeing-1}.
\heading{20.1  DSS de telemetría}{CALIDAD ANTES DE INTERPRETACIÓN}
El servicio consulta hasta 100 instantáneas desde el comienzo de la asignación ambiental activa hasta el reloj de consulta. El árbol separa dispositivos, descarta fechas futuras e IDs repetidos y exige estados válidos del firmware. La ventana de tendencias es de hasta 120 segundos; la última muestra caduca a los 30 segundos. Conservar historia no renueva la vigencia.
\begin{itemize}
\item SpO\textsubscript{2} menor de 90: prioridad de verificación. De 90 a menos de 95: revisión. Cero sin estado válido: dato no disponible.
\item ECG: con cinco valores válidos recientes como mínimo, compara la mediana de los dos más nuevos con la de los restantes. Una diferencia absoluta de al menos 20 lpm activa revisión.
\item Diferencia ECG/oxímetro de al menos 20 lpm: comprobar coherencia. Pulso óptico en el techo de formato de 300 lpm o salto de al menos 40 lpm en hasta 10 segundos: comprobar señal.
\item Respiración: informa valor y variación, sin clasificación clínica de normalidad. La radiación histórica no modifica el nivel fisiológico.
\end{itemize}
\begin{infobox}\textbf{Procedencia:} constantes demostrativas del código, no umbrales validados por NASA ni por ensayos clínicos del montaje. FDA y MedlinePlus se enlazan como orientación de comprobación, no como validación de estos parámetros.\end{infobox}
{\footnotesize\color{muted}Evidencia: \path{AssessmentService.ts}, \path{assessMeasurements.ts}, \path{telemetryDecisionTree.ts}, \path{TelemetryDecisionService.ts}, en \path{application/assessment} del backend.}
\newpage
\heading{21  Caso reproducible y trazabilidad}{DATOS SINTÉTICOS DE PRUEBA}
Ejemplo de \path{tests/telemetry-decision-tree.test.ts}, no un registro humano ni evidencia de precisión del sensor. Dos instantáneas sintéticas del mismo dispositivo y escenario, a 15 y 30 segundos antes del reloj evaluador:
\par\noindent\begin{tabularx}{\linewidth}{@{}lXX@{}}
\toprule Campo & Más reciente & Anterior\\\midrule
SpO\textsubscript{2} & 80 \% & 78 \%\\
Pulso óptico & 300 lpm & 17 lpm\\
Estado del oxímetro & 200 & 200\\
Pulso ECG / estado & 0 / 50 & 76 / 100\\
Respiración / estado & 0 / 50 & 17 / 100\\\bottomrule
\end{tabularx}\par
\begin{enumerate}[leftmargin=5mm,itemsep=5pt]
\item \textbf{Validación:} la última lectura es vigente. El estado óptico permite analizar SpO\textsubscript{2} y pulso. ECG y respiración actuales quedan no disponibles; sus ceros no son mediciones humanas.
\item \textbf{Evaluación:} resultado \texttt{priority} por SpO\textsubscript{2} de 80 \%. Pulso de 300 lpm: \texttt{review} y \texttt{checkRequired} por coincidir con el techo de formato.
\item \textbf{Acciones:} verificar oxigenación, contrastar pulso y recuperar ECG y respiración. La duda de calidad del pulso no elimina la prioridad de la saturación.
\item \textbf{Caducidad:} otra prueba con la última muestra a 31 segundos devuelve \texttt{insufficient}; retira advertencias fisiológicas actuales y solicita recuperar captura.
\end{enumerate}
El DSS de telemetría calcula resultados al consultar. Los avisos persistidos de mediciones personales siguen otra ruta: conservan ID de medición, regla, referencia, dirección y evidencia histórica. Admiten estados abierto, en seguimiento y cerrado. La agrupación visual no borra detecciones originales.
\heading{21.1  Seguridad y privacidad}{CONTROLES IMPLEMENTADOS}
La API exige sesión y permisos por rol. El propietario se resuelve desde la cuenta o vínculo del dispositivo; las consultas filtran por identidad y la captura valida el origen autorizado. Contraseñas y respuestas de recuperación usan hashes scrypt con sal. La cookie es HttpOnly, SameSite estricto y Secure en Vercel.

PostgreSQL/Supabase conserva cuentas, sesiones, mediciones, telemetría y seguimiento. El idioma y tema se guardan localmente. Cerrar sesión no elimina el historial; no se declara borrado automático ni certificación de privacidad. Estos controles no sustituyen una auditoría de seguridad o evaluación normativa.
\vfill
{\footnotesize\color{muted}Evidencia: \path{application/alerts}, \path{infrastructure/security/credentials.ts}, \path{transport/http.ts}; ingesta y permisos en \path{tests/cloud.test.ts}.}
\newpage
\heading{22  Verificación y resultados}{EJECUCIÓN DEL 15 NOVIEMBRE 2026}
Se ejecutó \texttt{npm test} sobre la base \texttt{b2daf00}, antes de ampliar este documento: \textbf{59 pruebas registradas, 58 aprobadas, 0 fallidas y 1 omitida}. La integración PostgreSQL se omitió por no disponer de \texttt{DATABASE\_URL} en esa ejecución; no acredita conexión real con Supabase.
\par\noindent\begin{tabularx}{\linewidth}{@{}>{\bfseries}p{35mm}X@{}}
\toprule Área & Cobertura observada\\\midrule
Captura BLE & Fragmentos, duplicados, once campos, UTF-8, desbordamiento, ausencia de contacto y estados.\\[6pt]
Evaluación personal & Calidad, caducidad, referencia separada, reposo confirmado y aislamiento de personas, entorno y simulación.\\[6pt]
DSS & Fechas futuras, intervalos ambientales, duplicados, dispositivos, prioridades, acciones y ausencia de influencia radiológica en el estado fisiológico.\\[6pt]
Presentación & Pulso histórico frente a actual, señal intermitente, segmentos ECG, filtro visual, grupos de avisos y eje temporal.\\[6pt]
Servicio DSS & Dependencias simuladas para permisos, asignación histórica y desconexión; no reemplazan integración real de base de datos.\\\bottomrule
\end{tabularx}\par
\textbf{Reproducción:} usar Node 22 compatible, instalar con \texttt{npm ci} y ejecutar \texttt{npm test}. La integración de \path{tests/cloud.test.ts} necesita una base de prueba autorizada y crea un esquema temporal. No publicar credenciales. Compilación y tipos son controles de software, no ensayos biomédicos.
\heading{22.1  Uso NASA y alcance de evidencia}{REFERENCIAS TRAZABLES}
El DSS importa \path{data/radlab_reference_summaries.json}. Son resúmenes suministrados: sin CSV originales no se reproduce limpieza, unidades de origen ni percentiles. CRaTER es orbital, MSL-RAD mide en Curiosity y EAD-MU01 en Orion; no son dosis del astronauta. Las fichas planetarias aportan valores documentales al catálogo. MEDA y Diviner no son flujos en vivo integrados.
\heading{22.2  Limitaciones y siguientes validaciones}{TRABAJO PENDIENTE}
\begin{itemize}
\item Documentar placa, pines, alimentación, colocación y versión real del firmware. No se añade un esquema eléctrico no verificado.
\item Comparar contra instrumentos de referencia y registrar error, condiciones y tamaño de muestra. No se acredita calibración ni validación clínica.
\item Probar captura física prolongada, movimiento, reconexión, pérdida de Internet y navegadores; registrar fallos y recuperación.
\item Recuperar CSV y metadatos originales para reproducir resúmenes; evaluar seguridad, privacidad y operación antes de uso fuera de demostración.
\end{itemize}
{\footnotesize\color{muted}Código y pruebas: \url{https://github.com/AMVMesias/ASTROVITAL}. Se distingue implementación, prueba automatizada y validación física pendiente.}
